\chapter{micro-ROS: cómo hablan la Pi y los ESP32}\label{cap:microros}

\section{Qué es micro-ROS}

Un ESP32 no puede correr ROS~2 completo (no tiene sistema operativo ni red DDS). \textbf{micro-ROS} es una versión
reducida que corre en el microcontrolador y habla con un \textbf{agente} en la Pi por un enlace simple (acá, el USB
serie). El agente traduce: para el resto de ROS~2, el ESP32 es un nodo más.

\begin{figure}[H]
\centering
\begin{tikzpicture}
  \node[placa,text width=34mm] (esp) at (0,0) {ESP32-S3\\\footnotesize firmware con micro-ROS\\(nodo \cod{base})};
  \node[caja,text width=36mm] (ag) at (6.2,0) {agente micro-ROS\\\footnotesize\cod{micro\_ros\_agent serial}};
  \node[nodo,text width=36mm] (ros) at (11.6,0) {resto de ROS 2\\\footnotesize\cod{mision}, \cod{teleop}, \cod{ros2 topic}\ldots};
  \draw[<->,thick,acento] (esp) -- node[topico,above=2pt]{USB serie} node[topico,below=2pt]{XRCE-DDS} (ag);
  \draw[<->,thick,acento] (ag) -- node[topico,above=2pt]{DDS (red)} (ros);
  \begin{scope}[on background layer]
    \node[grupo,fit=(ag)(ros),label={[font=\small\bfseries]above:Raspberry Pi 5}] {};
  \end{scope}
\end{tikzpicture}
\caption{El agente es el puente. Hay \textbf{un agente por placa}, cada uno en su puerto USB.}
\end{figure}

\subsection{Arrancar los agentes (en la Pi)}
\begin{lstlisting}[language=consola]
source /opt/ros/jazzy/setup.bash
source ~/microros_ws/install/local_setup.bash
# Base (ESP32-S3, USB nativo)
ros2 run micro_ros_agent micro_ros_agent serial --dev /dev/serial/by-id/usb-Espressif_USB_JTAG_serial_debug_unit_*-if00
# Brazo (WROOM con CH340) — a 115200 baudios: a 921600 el CH340 pierde datos en la Pi
ros2 run micro_ros_agent micro_ros_agent serial --dev /dev/serial/by-id/usb-1a86_USB_Serial-if00-port0
\end{lstlisting}

\begin{nota}[Flashear con el agente corriendo]
El agente tiene abierto el puerto USB: hay que detenerlo antes de flashear (\cod{pio run -e ros -t upload}) y volver a
arrancarlo después. Si la S3 no abre sesión, puede haber quedado en \textbf{modo descarga} (pasa si el puerto se abre con
DTR/RTS bajos). Se saca con:
\cod{\textasciitilde/venvs/esptool/bin/esptool --port <S3> --before no\_reset --after watchdog-reset read-mac}.
\end{nota}

\section{La sesión: conexión, hora y ping}

\begin{enumerate}
  \item Al arrancar, el ESP32 busca al agente. Mientras no lo encuentra, no publica nada.
  \item Al conectarse, crea su nodo, sus publicadores y suscriptores, y \textbf{sincroniza la hora} con el agente
        (todos los mensajes llevan la hora de la Pi). Sin hora, no publica.
  \item Cada \textbf{100~ms} hace un \textbf{ping} al agente. Si fallan 3 seguidos ($\sim$1,5~s con los reintentos),
        considera la sesión perdida: \textbf{detiene los motores apenas lo detecta} y queda \textbf{desarmado}.
  \item Vuelve a buscar al agente y, al reconectar, arranca una sesión nueva, otra vez desarmada.
\end{enumerate}

\cod{/<parte>/estado} incluye los contadores \cod{enlace: pings\_fallidos cortes}: sirven para ver si el USB anda
bien. En la prueba con la Pi a plena carga hubo 0 cortes.

\section{Armar y desarmar}\label{sec:armar}

\begin{idea}[Qué es armar]
Cada ESP32 tiene un \textbf{seguro por sesión}. Al empezar cada sesión con el agente queda \textbf{DESARMADO}: ignora
las órdenes de movimiento (\topic{/cmd\_vel} en la base; \topic{/brazo/joint\_cmd}, la pinza y los comandos que mueven
en el brazo). Para habilitarlo se publica el texto \cod{armar} en \topic{/base/comando} o \topic{/brazo/comando}.
\end{idea}

\begin{figure}[H]
\centering
\begin{tikzpicture}
  \node[estado] (sin) at (0,0) {SIN AGENTE};
  \node[estado] (des) at (5.5,0) {DESARMADO};
  \node[estado,draw=verde,fill=verde!8] (arm) at (11,0) {ARMADO};
  \draw[flecha] (sin) -- node[topico,above]{conecta + hora} (des);
  \draw[flecha] ([yshift=2mm]des.east) -- node[topico,above]{\cod{armar}} ([yshift=2mm]arm.west);
  \draw[flecha] ([yshift=-2mm]arm.west) -- node[topico,below]{\cod{desarmar}} ([yshift=-2mm]des.east);
  \draw[flecha,rojo] (arm.south) -- ++(0,-1.0) -| node[topico,pos=0.25,below,text=rojo]{3 pings perdidos ($\sim$1,5 s): frena los motores} (sin.south);
  \draw[flecha,rojo] (des.north) -- ++(0,0.9) -| node[topico,pos=0.25,above,text=rojo]{se pierde el agente} (sin.north);
\end{tikzpicture}
\caption{Estados de la sesión micro-ROS de cada ESP32.}
\label{fig:armar}
\end{figure}

\begin{itemize}
  \item \textbf{Desarmado acepta}: detener (\cod{s}), liberar (\cod{l}), borrar fallas (\cod{r}), cero (\cod{z}),
        diagnóstico (\cod{e}, \cod{sens}), \cod{cal}, \cod{wd}.
  \item \textbf{Desarmado rechaza}: \topic{/cmd\_vel}, \cod{v}, \cod{w}, \cod{d}, \cod{m}, \cod{j}, \cod{home},
        \topic{/brazo/joint\_cmd}, \cod{pinza}, \cod{giro}.
\end{itemize}

\subsection{Por qué existe}
Imaginá que el USB hace falso contacto y la placa se reconecta. Si un nodo seguía publicando ``avanzá a 0,1~m/s''
(por ejemplo, el teclado que quedó abierto), sin el seguro el robot arrancaría solo al volver la conexión. Con el
seguro, después de cada reconexión \textbf{alguien tiene que volver a decir \cod{armar} a propósito}.

\subsection{Quién arma}
\begin{itemize}
  \item A mano: \cod{ros2 topic pub --once /base/comando std\_msgs/msg/String} \cod{"\{data: 'armar'\}"}.
  \item \cod{mision} arma el brazo en su estado INICIO, reintentando cada segundo hasta ver \cod{ARMADO} en
        \topic{/brazo/estado}. No lo vuelve a armar después de una reconexión: eso queda como falla.
  \item \cod{recorrido} arma la base en su estado INICIO, igual.
  \item \cod{mision} \textbf{desarma} el brazo si ve \cod{SETA} en \topic{/base/estado}.
\end{itemize}

\subsection{Armar no es todo}
Estar armado solo significa ``acepto órdenes de esta sesión''. Las demás protecciones siguen activas: en la base, la
seta, el vuelco, los bumpers, la distancia y el watchdog; en el brazo, el cero, los límites, el atasco. Por ejemplo,
con la seta apretada la base puede estar armada pero no se mueve. Al soltar la seta hay que mandar \cod{r} (rearmar
la seguridad) además de que esté armada.

\section{Watchdog de comandos}

En la base, con micro-ROS el watchdog de comandos \textbf{arranca en 300~ms y no se puede apagar} (se puede cambiar
entre 50 y 1000~ms con \cod{wd}). Si \topic{/cmd\_vel} deja de llegar por más de ese tiempo, la base frena. Por eso
\cod{ros2 topic pub} para moverse se usa con \cod{-r 10} (10 veces por segundo) y no con \cod{--once}.

\begin{pendiente}
Medir con las ruedas en el piso la distancia real de frenado por el watchdog. Estimación (revisión de Codex): a
1~m/s, $\sim$30~cm antes de empezar a frenar; en la simulación, 11--13~cm desde 0,2~m/s.
\end{pendiente}

\section{Tópicos de la base}

{\small\begin{longtable}{@{}p{3.6cm} p{3.9cm} p{1.5cm} p{5.9cm}@{}}
\toprule
\textbf{Tópico} & \textbf{Tipo} & \textbf{Sentido} & \textbf{Detalle} \\
\midrule
\endhead
\topic{/odom} & \cod{nav\_msgs/Odometry} & sale, 25~Hz & \cod{odom} → \cod{base\_link} \\
\topic{/imu/data\_raw} & \cod{sensor\_msgs/Imu} & sale, 50~Hz & m/s² y rad/s, sin orientación \\
\topic{/range/frente}, \topic{/range/izq}, \topic{/range/der} & \cod{sensor\_msgs/Range} & sale, 10~Hz & cada sensor mide cada 120~ms; +Inf = nada al alcance; NaN = sin eco \\
\topic{/bumpers} & \cod{std\_msgs/UInt8} & sale, 10~Hz & bits 0..2: frente, izquierda, derecha \\
\topic{/battery} & \cod{sensor\_msgs/BatteryState} & sale, 1~Hz & tensión \\
\topic{/base/estado} & \cod{std\_msgs/String} & sale, 2~Hz & armado, modo, seguridad, bloqueos, IMU, fallas, enlace \\
\topic{/base/respuesta} & \cod{std\_msgs/String} & sale & lo que mostraría el monitor serie \\
\topic{/cmd\_vel} & \cod{geometry\_msgs/Twist} & entra & \cod{linear.x} [m/s], \cod{angular.z} [rad/s] \\
\topic{/base/comando} & \cod{std\_msgs/String} & entra & cualquier comando del monitor, más \cod{armar} y \cod{desarmar} \\
\bottomrule
\end{longtable}}

\section{Tópicos del brazo}

{\small\begin{longtable}{@{}p{3.6cm} p{3.9cm} p{1.5cm} p{5.9cm}@{}}
\toprule
\textbf{Tópico} & \textbf{Tipo} & \textbf{Sentido} & \textbf{Detalle} \\
\midrule
\endhead
\topic{/brazo/joint\_states} & \cod{sensor\_msgs/JointState} & sale, 20~Hz & \cod{j1..j3}: posición [rad], velocidad [rad/s] \\
\topic{/brazo/estado} & \cod{std\_msgs/String} & sale, 2~Hz & armado; modo, cero, duty, falla y \cod{llego} de cada articulación \\
\topic{/brazo/respuesta} & \cod{std\_msgs/String} & sale & lo que mostraría el monitor serie \\
\topic{/brazo/joint\_cmd} & \cod{sensor\_msgs/JointState} & entra & posiciones [rad] con nombres \cod{j1..j3} (o las 3 en orden). Se valida todo: se aplica todo o nada \\
\topic{/pinza/cierre} & \cod{std\_msgs/Float32} & entra & 0 = abierta, 100 = cerrada \\
\topic{/pinza/giro} & \cod{std\_msgs/Float32} & entra & grados, $-90$..$90$ \\
\topic{/brazo/comando} & \cod{std\_msgs/String} & entra & cualquier comando del monitor (\cod{z}, \cod{home 2}, \cod{s}, \cod{r}\ldots), más \cod{armar} y \cod{desarmar} \\
\bottomrule
\end{longtable}}

Si se pierde el agente, J1--J3 se detienen y el brazo queda desarmado; \textbf{los servos quedan donde están}, para
no soltar el limón.
